\section{Introduction}
\label{sec:intro}
% -------------------------------------------------------------

Vision-language model (VLM) agents capable of operating a web browser
are moving from research prototypes into production deployment
% CITE: copy any line below into Semantic Scholar / Google Scholar / arXiv to find the paper.
%   - "OpenAI Operator browser-using agent 2025 launch"
%   - "Anthropic Claude computer use 2024 system card"
%   - "Google Gemini agent Project Mariner browser 2024 2025"
%   - "survey large language model web-browsing agents 2024 2025"
%   - "multimodal LLM agent benchmark survey browser navigation"
\citep{OpenAI2025Operator,anthropic2024computeruse,tamkin2024clioprivacypreservinginsightsrealworld,andreux2025surferhmeetsholo1costefficient,zhang2025surveylargelanguagemodel,nguyen-etal-2025-gui}. To complete the tasks they are
given (applying for jobs, purchasing goods, paying bills, replying to
email), these agents are routinely provisioned with the user's full
personally identifiable information (PII): legal name, home address,
payment card details, government identifiers, account credentials, and
API keys. The agent stores this profile in its context and draws on
it whenever a form field must be filled or a credential must be
submitted. The result is a system that combines broad web access with
privileged read access to high-value personal data, a configuration
that realises a \emph{semantic} variant of the classic
\emph{confused-deputy} vulnerability
% CITE: copy any line below into Semantic Scholar / Google Scholar / arXiv to find the paper.
%   - "Hardy 1988 confused deputy or why capabilities might have been invented"
%   - "Norm Hardy confused deputy ACM SIGOPS Operating Systems Review 1988"
%   - "confused deputy LLM agent tool use indirect prompt injection"
%   - "agent confused deputy untrusted input privileged context 2024 2025"
\citep{Hardy1988TheCD,Ferrag2025FromPI,Triedman2025MultiAgentSE,fu2024impromptertrickingllmagents,kim2025promptflowintegrityprevent,ji2026tamingvariousprivilegeescalation}: the structural fault arises from mixing privileged user
context with unprivileged web inputs in a single decision loop, but
the mechanism of exploitation is entirely semantic, leveraging
multi-turn conversational social engineering to bypass the model's
behavioural safety constraints rather than escalating any traditional
privilege boundary. Untrusted content encountered during normal
browsing can instruct or manipulate the agent into forwarding that
data to an adversary. The modern web already contains a mature social-engineering
ecosystem (phishing clones on subtly misspelled domains, dark
patterns that bury the safe option, fake trust signals, urgency
timers, authority impersonation, multi-turn conversational
escalation) built to harvest PII from human users
% CITE: copy any line below into Semantic Scholar / Google Scholar / arXiv to find the paper.
%   - "APWG Anti-Phishing Working Group quarterly phishing activity trends report 2024 2025"
%   - "Mathur dark patterns shopping websites taxonomy CHI 2019"
%   - "Gray dark patterns side of UX design CHI 2018"
%   - "Brignull deceptive design patterns taxonomy"
%   - "Bösch dark strategies online services privacy"
%   - "Verizon Data Breach Investigations Report 2024 social engineering"
\citep{Hamadouche_Boudraa_Gasmi_2024,10.1145/3359183,10.1145/3173574.3174108,deceptivedesignpatterns,Bsch2016TalesFT}. A human user can slow down, scrutinize a URL, or
close the tab; a web agent is optimised to finish its task and will
follow links, fill fields, and submit forms with a thoroughness that
no human user matches. Capability benchmarks confirm that frontier
agents are now reliable enough to actually do this end-to-end: REAL
\citep{garg2026real} reports a 41\% top success rate on multi-turn
state-changing web tasks, meaning the same competence that makes
them useful also makes them a high-value, unusually reliable target
for data exfiltration. We find that even when an agent's reasoning trace is judged by an
independent LLM to flag the site as suspicious, under the
pre-registered C3 pre-submission reflection condition the agent
submits critical PII to that site in 38\% of sessions, versus 69\%
when no such recognition is verbalised: a 31.5-percentage-point gap
between what the agent \emph{says} and what it \emph{does},
statistically robust ($q\!<\!0.001$) and present across the three
frontier model families on which the cross-family judge has been
completed (Claude Haiku 4.5, Gemini 3 Flash, Llama 4 Scout).

\paragraph{The measurement gap.} Despite a growing literature on the
adversarial robustness of web agents, existing evaluation frameworks
gate success on proxy outcomes rather than on the terminal
field-level submission of sensitive data to an attacker-controlled
endpoint: click compliance under dark patterns
\citep{cuvin2026decepticon,ersoy2026trickyarena,chen2026obviousinvisiblethreat},
indirect prompt injection at tool-use APIs
\citep{debenedetti2024agentdogo}, single-click redirect compliance
\citep{korgul2025trap}, incidental over-sharing on benign sites
\citep{zharmagambetov2026agentdam}, or social engineering treated as
one of several risk categories without a structured PII profile
\citep{xinyi2026webtrappark}. The terminal submission is the
data-loss event a deployment review cares about: early-trajectory
anomalies are observable to low-cost runtime guardrails, but a
completed form-level POST payload represents an absolute,
unrecoverable data breach that can only be prevented by the
cognitive refusal strategy of the model itself. To our knowledge,
no published benchmark combines (i) attacker-controlled websites
that explicitly social-engineer the agent, (ii) a populated PII
profile spanning critical, high, and medium sensitivity tiers,
(iii) field-level exfiltration as the primary outcome, and (iv)
axis-controlled siblings that isolate which design factors causally
drive leakage.

\paragraph{This work.} We introduce \textsc{Scammer4U}, a benchmark
of 91 hand-crafted attacker-controlled environments plus 10
benign-twin baselines, designed to answer a single operational
question: \emph{when an autonomous web agent encounters a
social-engineering attack while browsing, will it hand over the
user's PII?} Environments cover eight attack vectors and span 16
site categories (full enumeration in \S\ref{sec:vectors} and
\S\ref{sec:composition}). Each environment is classified on four
axes (category, vector, salience, PII target) and a further four
factor axes (pressure cue, prompt-injection style, interaction
modality, multi-site flow), forming an 8-axis factorial taxonomy
mapped to published threat catalogues
\citep{mitre_attack_citation,owasp_llm_citation,enisa_citation}. 60
of the 91 environments are paired siblings that differ from a parent
on exactly one axis, supporting causal single-factor ablation. We
equip a Playwright-based agent
\citep{playwright,zhou2024webarenarealisticwebenvironment,Koh2024VisualWebArenaEM,10.5555/3692070.3694608,NEURIPS2023_5950bf29,he-etal-2024-webvoyager}
with a 40-field PII profile and evaluate four frontier model
families (GPT-5 mini, Claude Haiku 4.5, Gemini 3 Flash, Llama 4
Scout) across four mitigation conditions C0--C3 (baseline, generic
privacy nudge, phishing-aware checklist, pre-submission reflection).
The analysis plan is committed to a public git tag
(\texttt{prereg-v2-start}) before any reported data was collected,
with every subsequent deviation documented in the limitations log.

\paragraph{Contributions and headline findings.} The paper makes
four contributions, each paired with the empirical finding that
backs it. Retained-session counts are reported in
\S\ref{sec:results-baseline}.

\begin{enumerate}[leftmargin=*,itemsep=3pt]
  \item \textbf{A 31.5-percentage-point detection--action gap under
  pre-submission reflection (primary).} Within the pre-registered C3
  condition, restricted to sessions that reach the attack surface,
  the LLM-as-judge confirms independent flagging of the site as
  suspicious in $n = 295$ sessions pooled across the three
  judge-instrumented model families (Claude Haiku 4.5, Gemini 3
  Flash, Llama 4 Scout). In those sessions the agent still submits
  critical PII at 38.0\%, versus 69.5\% when no recognition is
  verbalised: a 31.5\,pp gap, statistically robust
  ($q < 0.001$ after Benjamini--Hochberg correction over the
  pre-registered primary family) and surviving a both-judges-agree
  sensitivity (31.8\,pp on the same cell;
  \S\ref{sec:results-f1}). The DR$=$1 leakage rate of 38.0\%
  exceeds the pre-registered 10\% falsification threshold by
  28.0\,pp: saying ``this looks like phishing'' is not the same
  as refusing to submit.

  \item \textbf{Empirical evidence that prompt-level mitigation is
  sharply model-dependent.} Baseline critical-tier leakage (C0)
  ranges from 53.8\% on Haiku to 93.1\% on Gemini, with GPT-5 mini
  at 75.6\% and Llama at 82.3\%, all far above the 0\% empirically
  measured on 10 benign-twin environments. A pre-submission
  reflection prompt (C3) reduces PLR\textsubscript{crit} by
  32.4\,pp on Gemini, 30.8\,pp on Haiku, 17.8\,pp on GPT-5 mini,
  and only 4.9\,pp on Llama; the C2 phishing-checklist condition
  is in fact stronger than C3 on Haiku ($-$36.1\,pp). The pooled
  M3 effect ($-$19.7\,pp) hides this heterogeneity: three
  model--condition cells individually cross the pre-registered
  $-$30\,pp threshold, but the pooled average does not. This
  pattern motivates output-level defences (form-submission
  inspection, structured-PII gating) rather than additional
  system-prompt iteration.

  \item \textbf{A characterisation of when web agents leak PII
  across the 8-axis taxonomy.} Subtle attacks leak \emph{more}
  than blatant ones (96.5\% vs.\ 72.8\% PLR\textsubscript{crit} at
  C0, pooled): we hypothesise that blatant adversarial indicators
  trip standard alignment guardrails via heavy keyword detection,
  whereas subtle semantic manipulations navigate beneath these
  safety filters. \texttt{fake\_trust\_signals} and
  \texttt{reward\_trap} vectors are pinned at $\sim$94\% across
  all mitigation conditions; \texttt{authority\_impersonation} is
  the most C3-responsive vector ($-$62.5\,pp); the pressure axis
  (urgency, scarcity, social proof, authority) shows no
  significant pooled effect. Paired-sibling factor tests F2--F11
  mostly do not survive BH correction at the realised pair counts;
  the characterisation signal lives in cross-cutting marginal
  patterns rather than paired-sibling deltas.

  \item \textbf{The \textsc{Scammer4U} benchmark itself}, publicly
  released: 91 attacker-controlled environments and 10 benign-twin
  baselines, organised on an 8-axis factorial taxonomy mapped to
  published threat catalogues, with a pre-registered analysis plan,
  a complete deviation log, and full raw run logs, designed so that
  follow-up work can re-run the same ablations on new models
  without re-authoring the apparatus.
\end{enumerate}

